\speciestitle{Nitric Acid}{HNO\cs3}{HNO\cs3}
\begin{description}
\item[Useful range:] 215\,--\,1.5\,hPa (1.0\,hPa under enhanced conditions)
\item[Contact:] Gloria Manney, \email{manney@nwra.com}
\end{description}

\subsection*{Introduction}

The quality and reliability of the Aura MLS v2.2 HNO\cs3 measurements
were assessed in detail by \citet{SanteeEtal07hno3}.  The HNO\cs3 in
\thisv has been greatly improved over that in version v2.2; in particular,
a low bias through much of the stratosphere (especially evident at
levels with pressure greater than or equal to 100~hPa) has been largely
eliminated.  Figure~\ref{fig:zmv2v3diffs} shows an example of typical
differences between v2.2 and \thisv HNO\cs3\@.  Improvement in HNO\cs3
resulted from indirect effects of adding interline interference terms to
the O\cs3 line shape model, an updated CO line width parameter, using a
different 240 GHz channel configuration for retrieving HNO\cs3, and a
change the manner in which continuum signals are accounted for (see
\ref{sec:V3vsV2Overview}); these changes contribute approximately
equally to the HNO\cs3 improvement.  However, an unfortunate side effect
of the change to continuum handling is that it is more adversely
affected by clouds, causing spikes in the retrieval of 240-GHZz products
including HNO\cs3.  In addition, it also appears that the new continuum
treatment has led to a noisier HNO\cs3 product in the UTLS than that in
v2.2.  Lower \thisv values in Figure~\ref{fig:zmv2v3diffs} in the tropics
at the lowest levels result largely from these effects.

The MLS \thisv HNO\cs3 data are scientifically useful over the range 215
to 1.5\,hPa; values at 1\,hPa are also expected to be scientifically
useful under conditions of enhanced HNO\cs3 in the upper stratosphere,
but should be used with caution and in consultation with the MLS team.
HNO\cs3 values in the upper stratosphere, at 3.2 through 1.0\,hPa, are
frequently very low and may require averaging (this will usually be the
case at 1.5 and 1\,hPa, where the values are also noisier than at lower
levels), but during periods of enhancement in the upper stratosphere,
coherently evolving atmospheric signals with realistic morphology are
seen in individual daily maps.  The standard HNO\cs3 product is derived
from the \mbox{240-GHz} retrievals at pressures equal to or greater than
22\,hPa and from the \mbox{190-GHz} retrievals for lesser pressures.
The \texttt{Quality} and \texttt{Convergence} information included in
the standard HNO\cs3 files are from the \mbox{240-GHz} retrievals, and
\emph{apply only to pressures 22\,hPa or greater} (see the data
screening discussion below).

A summary of the precision and resolution (vertical and horizontal) of
the \thisv HNO\cs3 measurements as a function of altitude is given in
Table~\ref{tab:HNO3Summary}.  The impact of various sources of
systematic uncertainty was quantified for v2.2, and it is expected that
these estimates will be similar for \thisv.  Table~\ref{tab:HNO3Summary}
also includes estimates of the potential biases and scaling errors in
the measurements compiled from the v2.2 uncertainty analysis (to be
updated for \thisv at a later date).  The overall uncertainty for an
individual data point is determined by taking the root sum square (RSS)
of the precision, bias, and scaling error terms (for averages, the
single-profile precision value is divided by the square root of the
number of profiles contributing to the average).  More details on the
precision, resolution, and accuracy of the MLS \thisv HNO\cs3 measurements
are given below.

\begin{figure}
\dontincludegraphics[angle=90,width=\linewidth]%
{figures/InspectZMContour_HNO3--v2-2_HNO3--v3-3_2005m08_LatPress-All}
\caption{V2.2 (top left) and v3.3 (top right) zonal mean HNO\cs3 for
  August~2005, and differences (v3.3\,$-$\,v2.2) expressed in ppbv (bottom
  left) and percent (bottom right).}
\label{fig:zmv2v3diffs}
\end{figure}


\subsection*{Resolution}
\onestandardavk{HNO3}{HNO\cs3}%

The resolution of the retrieved data can be described using `averaging
kernels' \citep[e.g.,][]{Rodgers00}; the two-dimensional nature of the
MLS data processing system means that the kernels describe both vertical
and horizontal resolution.  Smoothing, imposed on the retrieval system
in both the vertical and horizontal directions to enhance retrieval
stability and precision, reduces the inherent resolution of the
measurements.  Consequently, the vertical resolution of the \thisv HNO\cs3
data, as determined from the full width at half maximum of the rows of
the averaging kernel matrix shown in Figure~\ref{fig:AvkHNO3}, is
3\,--\,4\,km through most of the useful range, degrading to \texttildelow5\,km
at 22\,hPa and some levels in the upper stratosphere (see
Table~\ref{tab:HNO3Summary}).  Note that the averaging kernels for the
215 and 316\,hPa retrieval surfaces overlap over most of their depth,
indicating that the 316\,hPa retrieval provides little independent
information.  Figure~\ref{fig:AvkHNO3} also shows horizontal averaging
kernels, from which the along-track horizontal resolution is determined
to be 450\,--\,500\,km over most of the vertical range, improving to
250\,--\,300\,km between 15 and 4.6\,hPa, and degrading to
600\,--\,750\,km at 1.5 and 1\,hPa.  The cross-track resolution, set by
the widths of the fields of view of the \mbox{190-GHz} and
\mbox{240-GHz} radiometers, is \texttildelow10\,km.  The along-track separation
between adjacent retrieved profiles is 1.5\degsym\ great circle angle
(\texttildelow165\,km), whereas the longitudinal separation of MLS
measurements, set by the Aura orbit, is 10\degsym\,--\,20\degsym\ over
low and middle latitudes, with much finer sampling in the polar regions.

\subsection*{Precision}

\begin{figure}
\centering\dontincludegraphics[width=8.4cm]{figures/ComparePrecision2PanelHNO3}
\caption{Precision of the (left) v3.3 and (right) v2.2 MLS HNO\cs3
  measurements for four representative days (see legend).  Solid lines
  depict the observed scatter in a narrow equatorial band (see text);
  dotted lines depict the theoretical precision estimated by the
  retrieval algorithm.}
\label{fig:PrecProf}
\end{figure}

The precision of the MLS HNO\cs3 measurements is estimated empirically
by computing the standard deviation of the profiles in the
20\degsym-wide latitude band centered around the equator, where natural
atmospheric variability should be small relative to the measurement
noise.  Because meteorological variation is never completely negligible,
however, this procedure produces a upper limit on the precision-related
variability.  As shown in Figure~\ref{fig:PrecProf}, the observed
scatter in the \thisv data is \texttildelow0.6\,--\,0.7\,ppbv throughout the range
from 100 to 3.2\,hPa, below and above which it increases sharply.  The
scatter is essentially invariant with time, as seen by comparing the
results for the different days shown in Figure~\ref{fig:PrecProf}.

The single-profile precision estimates cited here are, to first order,
independent of latitude and season, but it should be borne in mind that
the large geographic variations in HNO\cs3 abundances gives rise to wide
range `signal to noise' ratios.  At some latitudes and altitudes and in
some seasons, HNO\cs3 abundances are smaller than the single-profile
precision, necessitating the use of averages for scientific studies.  In
most cases, precision can be improved by averaging, with the precision
of an average of $N$ profiles being 1/$\sqrt{N}$ times the precision of
an individual profile (note that this is not the case for averages of
successive along-track profiles, which are not completely independent
because of horizontal smearing).

The observational determination of the precision is compared in
Figure~\ref{fig:PrecProf} to the theoretical precision values reported
by the Level~2 data processing algorithms.  Although the two estimates
compare very well between 100 and 32\,hPa, above 22\,hPa the predicted
precision substantially exceeds the observed scatter.  This indicates
that the a priori information and the vertical smoothing applied to
stabilize the retrieval are influencing the results at the higher
retrieval levels.  Because the theoretical precisions take into account
occasional variations in instrument performance, the best estimate of
the precision of an individual data point is the value quoted for that
point in the L2GP files, but it should be borne in mind that this
approach overestimates the actual measurement noise at pressures less
than 22\,hPa.  Conversely, the observed scatter at pressures higher than
100\,hPa is considerably larger than the theoretical precision.  This is
related to the spikes and increased noise in the UTLS in \thisv versus
v2.2 HNO\cs3 mentioned above.  Procedures for screening outliers in this
region are discussed below.

\subsection*{Accuracy}

The effects of various sources of systematic uncertainty (e.g.,
instrumental issues, spectroscopic uncertainty, and approximations in
the retrieval formulation and implementation) on the MLS v2.2 HNO\cs3
measurements were quantified through a comprehensive set of retrieval
simulations; results for \thisv, to be completed at a later date, are
expected to be similar.  The results of the v2.2 uncertainty analysis
are summarized in Table~\ref{tab:HNO3Summary}; see
\citet{SanteeEtal07hno3} for further details of how the analysis was
conducted and the magnitude of the expected biases, additional scatter,
and possible scaling errors each source of uncertainty may introduce
into the data.  In aggregate, systematic uncertainties are estimated to
induce in the HNO\cs3 measurements biases that vary with altitude
between $\pm$0.5 and $\pm$2\,ppbv and multiplicative errors of
$\pm$5\,--\,15\% through most of the stratosphere, rising to
\mbox{\texttildelow$\pm$30\%} at 215\,hPa and \texttildelow50\% at and above 2.2\,hPa.
These uncertainty estimates are generally consistent with the results of
comparisons with correlative datasets, as discussed briefly below.

\subsection*{Data screening -- all data}
\label{sec:HNO3screening}

\begin{description}
\item[Pressure range:] \screeningrule{215\,--\,1.5\,hPa}
  \standardrangestatement
\standardprecisionrule
\evenstatusrule
\end{description}
\subsection*{Data screening -- upper troposphere, lower stratosphere
  (pressures of 22\,hPa or greater)}

The \texttt{Quality} and \texttt{Convergence} fields included in the
standard HNO\cs3 files are appropriate for use in screening at levels at
and below (that is, pressures greater than) 22\,hPa.  For those levels:

\begin{description}
\item[Quality:] \qualityrule{0.5}

  This threshold for \texttt{Quality} typically excludes \texttildelow2\,--\,4\% of
  HNO\cs3 profiles on a daily basis; it is a conservative value that
  potentially discards a significant fraction of ``good'' data points
  while not necessarily identifying all ``bad'' ones.

\item[Convergence:] \convergencerule{1.4}

  On a typical day this threshold for \texttt{Convergence} discards a
  very small fraction of the data, but on occasion it leads to the
  elimination of \texttildelow0.5\,--\,1\% of the HNO\cs3 profiles.
\item[Clouds:] \screeningrule{Clouds impact HNO\cs3 data in the UTLS,
    see discussion below and the discussion on `outliers' that follows.}

  Nonzero but even values of \texttt{Status} indicate that the profile
  has been marked as questionable, typically because the measurements
  may have been affected by the presence of thick clouds.  Globally
  \texttildelow10\,--\,15\% of profiles are identified in this manner, with the
  fraction of profiles possibly impacted by clouds rising to
  \texttildelow25\,--\,35\% on average in the tropics.  Clouds generally have
  little influence on the stratospheric HNO\cs3 data.  In the lowermost
  stratosphere and upper troposphere, however, thick clouds can lead to
  spikes in the HNO\cs3 mixing ratios in the equatorial regions.

  Therefore, it is recommended that at and below 68\,hPa, in addition to
  performing the `outlier screening' described below, all profiles with
  nonzero values of Status be discarded because of the potential for
  cloud contamination.  While this will reject some profiles that are
  probably not significantly impacted by cloud effects, it does
  eliminate a significant number of unphysical profiles that are not
  flagged by the outlier screening.

\item[Outliers:] \screeningrule{Alternative screening approaches in the
    UTLS remove outliers while reducing `false positives'}

  Outliers in \thisv HNO\cs3 at levels between 316 and 100 (sometimes to
  68) hPa frequently appear as highly negative mixing ratios at the
  lowest several retrieval levels, often as part of oscillatory profiles
  with unrealistically high values at higher altitudes.  A simple
  procedure is recommended to screen such profiles based on eliminating
  all profiles with large negative mixing ratios at pressure levels
  between 316 and 68\,hPa.  Through extensive examination of data
  screened in this way, flagging profiles that have either HNO\cs3 vmr
  less than $-$2.0\,ppbv at 316\,hPa or less than $-$1.6\,ppbv at any
  level between 215 and 68\,hPa eliminates most of the troublesome
  outliers, including those with positive vmr spikes overlying the
  negative ones that are directly flagged by these criteria.  This
  screening procedure is recommended for any studies focusing on the
  UTLS\@. That it effectively removes most of the suspect profiles that
  are not eliminated by requiring Status to be zero was evaluated as
  described in the following paragraphs:

  We compared the outlier screening method described above with a
  procedure based on using MLS cloud information: If the MLS ice water
  content (IWC) at 147\,hPa is greater than 0.003\,g/m$^3$ (indicating
  the presence of some cloud) for a profile, that profile, and the
  immediately adjacent ones along the orbit track are flagged (adjacent
  profiles are flagged assuming that the 1-D nature of the IWC retrieval
  versus the 2-D nature of the HNO\cs3 retrieval results in some
  uncertainty in the relative location of the cloud signal with respect
  to the trace gas profile).  Figure~\ref{fig:outlierfraction} shows the
  fraction of points eliminated by this procedure and the simpler
  recommended procedure based on direct identification of unphysical
  mixing ratios; the simpler procedure that is recommended compares very
  favorably with the more rigorous procedure based on cloud information:
  Each procedure eliminates a similar fraction of profiles (4\%
  globally, 15\% in the tropics), and is effective for removing most of
  the outliers.  Screening using MLS IWC as described here could be used
  (or compared with the recommended procedure) in analyses that are
  expected to be especially sensitive to the exact values in the
  tropical UTLS, but is not needed to obtain a high-quality HNO\cs3
  dataset in most cases.

  Figure~\ref{fig:outlierprofs} shows an example of the results of
  screening profiles by each of the \texttt{Quality},
  \texttt{Conver\-gence} and the recommended outlier flagging on a typical
  `bad' day (i.e., one with a relatively large number of outliers).  The
  \texttt{Quality} screening removes many of the profiles that are
  strongly negative at the bottom, and most or all of the remainder of
  these are flagged by the outlier screening; many of these profiles are
  oscillatory, so this screening also removes most or all of the strong
  positive outliers (typically at 147\,hPa).  Most of the profiles
  flagged by any of the criteria are in the tropics, as expected; the
  figure indicates that, at southern hemisphere high latitudes, low
  values of HNO\cs3 associated with the denitrified polar vortex are not
  triggering the outlier flagging.  As is often (but not always) the
  case, it is not clear in this example that the profiles flagged only
  by \texttt{Convergence} are unphysical or extreme.
 
\end{description}

\subsection*{Data screening -- upper stratosphere (pressures of 15\,hPa
  or less)}

The above screening criteria {\it should not be used} for 15\,hPa and
higher altitudes, as they result in filtering profiles for which all
quality indicators are good when the \texttt{Quality} and
\texttt{Convergence} values are properly taken from the 190-GHz HNO\cs3
information, and not filtering ones with indications of poor quality.
For any studies focusing on the upper stratosphere, it is highly
recommended that the user read from the \texttt{L2GP-DGG} files to
obtain the appropriate \texttt{Quality} and \texttt{Convergence} values
for the 190-GHz HNO\cs3 (from the \texttt{HNO3-190} swath), and use them
to apply the following screening criteria:
 
\begin{description}
\item[Clouds:] \screeningrule{Profiles where the \texttt{Status} field
    for \texttt{HNO3-190} has a non-zero even number can be used without
    restriction.} Clouds generally have little influence on the
  stratospheric HNO\cs3 data at these altitudes.

\item[Quality:] \screeningrule{Only profiles with a value of the
    \texttt{Quality} field for \texttt{HNO3-190} (see
    section~\ref{sec:StatusQuality}) \emph{greater} than 1.0 should be
    used in scientific study.}

  This threshold for \texttt{Quality} typically excludes
  \texttildelow1\,--\,3\% of HNO\cs3 profiles on a daily basis; it is a
  conservative value that potentially discards a significant fraction of
  ``good'' data points while not necessarily identifying all ``bad''
  ones.

\item[Convergence:] \screeningrule{Only profiles with a value of
    the \texttt{Convergence} field (see section~\ref{sec:StatusQuality})
    for the \texttt{HNO3-190} product \emph{less} than 1.6 should be
    used in investigations.}

  On a typical day this threshold for \texttt{Convergence} discards
  \texttildelow0.5\,--\,1.5\% of the HNO\cs3 profiles.

\item[Outliers:] For levels at and above (pressures less than) 4.6\,hPa,
  especially at 2.2\,hPa and above, some profiles show vertically
  oscillatory behavior in conditions where HNO\cs3 is very low.  The
  \texttt{Quality} and \texttt{Convergence} criteria defined above, when
  used together, eliminate many of these profiles; screening using both
  of these thresholds is thus particularly important.
\end{description}

\begin{figure}
\centering\dontincludegraphics[width=8.4cm]{figures/HNO3_frac_screened_by_lat_2}
\caption{Fraction of profiles flagged by two suggested outlier screening
  procedures for the UTLS as a function of latitude for \thisv data
  during 2006.}
\label{fig:outlierfraction}
\end{figure}

\begin{figure}
\centering\dontincludegraphics[width=16.4cm]{figures/profqual_utls2005d230}
\caption{HNO\cs3 profiles on 18~Aug~2005 color-coded by screening.
  Cyan profiles have \texttt{Quality} less than 0.5, olive-green
  \texttt{Convergence} greater than 1.4, and red both \texttt{Quality}
  less than 0.5 and \texttt{Convergence} greater than 1.4.  Orange
  profiles are those flagged by the simple screening procedure described
  above (using large negative mixing ratios at high pressures) after the
  profiles that failed \texttt{Quality} and/or \texttt{Convergence}
  tests were removed.  Black profiles are all those remaining (the
  `good' profiles) after the screening.  The left panels show all
  individual profiles in the day; the right panels show the means in
  each category, with the standard deviation shown as bars and the range
  as dotted lines.  The horizontal line is at 22\,hPa, above which
  HNO\cs3 is from the 190-GHz radiometer and thus not appropriately
  screened by these criteria.  The four pairs of panels show all
  profiles (top left), profiles between $-$20 and 20\degsym\ latitude (top
  right), profiles between $-$90 and $-$50\degsym\ latitude (bottom left) and
  profiles between 50 and 90\degsym\ latitude.}
\label{fig:outlierprofs}
\end{figure}

\subsection*{Review of comparisons with other datasets}


V3.3 / v3.4 HNO\cs3 has been compared with correlative datasets from a
variety of different platforms.  A consistent picture has emerged of
much closer agreement in \thisv than v2.2 with HNO\cs3 measurements from
ground-based, balloon-borne, and satellite instruments, especially in
the upper troposphere through the mid-stratosphere where MLS v2.2
HNO\cs3 mixing ratios were uniformly low by 10\,--\,30\%.  Example
comparisons with balloon-borne measurements
(Figures~\ref{fig:ftsumnerhno3comps} and \ref{fig:gbmshno3comps}) and
Atmospheric Chemistry Experiment-Fourier Transform Spectrometer
satellite measurements (Figure~\ref{fig:acehno3comps}) are shown.  The
GBMS ground-based measurements (Figure~\ref{fig:gbmshno3comps})
highlight the change from v2.2 to \thisv; in all years, MLS HNO\cs3
values increased from v2.2 to \thisv over most/all of the altitude
range, and in 2004\,--\,2005 and 2006\,--\,2007 show much closer
agreement with GBMS measurements; in all years, MLS HNO\cs3 agrees with
GBMS within the error bars.  The 2005\,--\,2006 winter was characterized
by extremely strong dynamical activity, which may contribute to the
different relationship between MLS and GBMS measurements in that year;
detailed GBMS/MLS comparisons are described by \citet{FiorucciEtal13}.

The Ft.\ Sumner balloon comparisons also show much improved agreement
between HNO\cs3 measured by several instruments with \thisv MLS data
(compare Figure~\ref{fig:ftsumnerhno3comps} with Figures~11 and 12 of
\citet{SanteeEtal07hno3}).  ACE-FTS comparisons also show improvements,
with the low bias in MLS virtually eliminated over the entire altitude
range shown in the top panel comparing with ACE-FTS v2.2 data in
May~2008 (other months show similar results) -- this can be contrasted
with Figure~25 of \citet{SanteeEtal07hno3}, which showed a low bias in
MLS v2.2 data with respect to ACE v2.2 throughout the useful altitude
range.

ACE-FTS HNO\cs3 reprocessed with v3.0 shows improvements in the UTLS,
and is useful up to \texttildelow60\,km. The bottom panels of
Figure~\ref{fig:acehno3comps} show a comparison of ACE-FTS v3.0 data with MLS \thisv
data during January 2005, indicating good agreement throughout the
altitude range.

Preliminary comparisons also indicate closer agreement for version 3
HNO\cs3 with aircraft measurements in the UTLS than for version 2.

\begin{figure}
\centering\dontincludegraphics[width=13.4cm]{figures/ftsumnerallhno3}
\caption{Comparisons with balloon-borne measurments at Ft. Sumner in
  2004 (left) and 2005 (right).  (Top panels) Path traversed by
  measurements from the balloon-borne MkIV (blue triangles) and FIRS-2
  (green and orange crosses represent two separate profiles) instruments
  during the flights from Ft.\ Sumner, NM, on 23\,--\,24~September 2004
  (left) and 20\,--\,21~September~2005 (right).  Measurement tracks from
  nearby MLS orbits are also shown (open circles).  The two MLS data
  points closest to the balloon measurements in time and space are
  indicated by red squares, with the closer one denoted by a filled
  symbol; the 500-km radius around the closest MLS point is overlaid in
  black.  (Bottom) Profiles of HNO\cs3 from MLS (red squares), MkIV
  (blue triangles), and FIRS-2 (green and orange crosses), corresponding
  to the symbols in the top panel.  Error bars represent the estimated
  precisions of each instrument, taken from the data files.  }
\label{fig:ftsumnerhno3comps}
\end{figure}

\begin{figure}
\vspace*{-0.5cm}
\centering\dontincludegraphics[width=13.4cm]{figures/gbmsallhno3}
\caption{Comparisons with GBMS balloon measurements.  (Top) Averages of
  all GBMS (blue) profiles and closest MLS (v2.2 in red, \thisv in
  green) coincidences at Testa Grigia (45.9$\degsym$N, 7.7$\degsym$E)
  during the 2004\,--\,2005, 2005\,--\,2006, and 2006\,--\,2007 winters;
  bars are standard deviation of the mean.  (Center) Differences between
  mean profiles shown in top panels (GBMS\,--\,MLS) in ppbv; (bottom)
  same differences, expressed as percentages.  }
\label{fig:gbmshno3comps}
\end{figure}

\begin{figure}
\centering\dontincludegraphics[width=15.4cm]{figures/acehno3compsv2v3}
\caption{Comparisons with ACE-FTS v2.2 measurements during May~2008
  (top) and ACE-FTS v3.30 measurements during Jan~2005 (bottom).  (left)
  Global ensemble mean profiles of the collocated matches for both
  instruments (MLS, red; ACE-FTS, blue). (middle) Mean percentage
  difference profiles between the two measurements (MLS\,--\,ACE-FTS)
  (cyan); standard deviation about the mean differences (Observed SD;
  orange) and the percentage root sum square of the precisions on both
  instrument measurements (Expected SD; magenta).  (right) As in Figure
  13 (middle) except plotted in mixing ratio.}
\label{fig:acehno3comps}
\end{figure}

\subsection*{Desired improvements for future data version(s)}
\begin{itemize}
\item Reduce noise/spikes in UTLS and in upper stratosphere.
\item Minimize the impact of thick clouds on the retrievals to further
  improve the HNO\cs3 measurements in the upper troposphere and
  lowermost stratosphere.
\end{itemize}

\begin{table*}
\caption{Summary of Aura MLS \thisv HNO\cs3 Characteristics}\vspace{0.5ex}
\label{tab:HNO3Summary}
\begin{minipage}{\linewidth}
  \renewcommand{\thefootnote}{\thempfootnote}
  \centering\begin{tabular}{cccccl} \toprule
\parbox[m]{18mm}{\bfseries\centering{Pressure\\ / hPa}} &
\parbox[m]{24mm}{\bfseries\centering{Resolution\\ \mbox{V $\times$ H}%
\footnote{Horizontal resolution in along-track direction.} \\ / km}} &
\parbox[m]{18mm}{\bfseries\centering{Precision%
\footnote{Precision on individual profiles, determined from observed
 scatter in the data in a region of minimal atmospheric variability.} \\/ ppbv}} &
\parbox[m]{20mm}{\bfseries\centering{Bias\\uncertainty%
\footnote{Values should be interpreted as \mbox{2-$\sigma$} estimates
 of the probable magnitude.}\\ / ppbv}} &
\parbox[m]{20mm}{\bfseries\centering{Scaling\\uncertainty%
\footnotemark[\value{mpfootnote}]\\/ \%}} &
\textbf{Comments} \\
%
\midrule
%
0.68\,--\,0.001 & ---  & ---    & ---    & --- & Unsuitable for scientific
  use\\
%
1.0 & 4 $\times$ 650\,--\,750  & $\pm$1.2 & $\pm$0.5 & $\pm$50\% & Caution, averaging
recommended \\
%
1.5 & 4.5 $\times$ 550\,--\,600 & $\pm$1.0  & $\pm$0.5 & $\pm$50\% & Averaging recommended \\
%
2.1 & 5 $\times$ 500 & $\pm$0.9  & $\pm$1.0 & $\pm$50\% & \\
%
3.2 & 4.5 $\times$ 400 & $\pm$0.7  & $\pm$0.5  & $\pm$10\,--\,15\% & \\
%
15\,--\,6.8  & 3\,--\,4 $\times$ 250\,--\,300 & $\pm$0.7 & $\pm$1\,--\,2  &
  $\pm$10\% & \\
%
22 & 5 $\times$ 450 & $\pm$0.7 & $\pm$1\,--\,2 &
$\pm$10\% & \\
%
100\,--\,32  & 3\,--\,4 $\times$ 350\,--\,400 & $\pm$0.7 & $\pm$0.5\,--\,1 &
$\pm$5\,--\,10\% & \\
%
147        & 3.5  $\times$ 400\,--\,450      & $\pm$0.8 &
$\pm$0.5  & $\pm$15\% & \\
%
215        & 3.5\,--\,4 $\times$ 500      & $\pm$1.2 & $\pm$1    &
\texttildelow$\pm$30\% & \\
%
316        & ---      & ---    & ---    & --- & Unsuitable for
scientific use\\
%
1000\,--\,464  & ---      & ---    & ---    & --- & Not retrieved\\
%
\bottomrule
\end{tabular}\par
\end{minipage}
\end{table*}

%%% Local Variables: 
%%% mode: latex
%%% TeX-master: "v4-x-report"
%%% End: 

